% =============================================================================
%  Remerciements
%
%  À relire et à ajuster avant remise : les remerciements doivent être ceux de
%  l'auteur. Les noms cités ici sont ceux consignés dans le dossier de projet.
% =============================================================================

\chapter*{Remerciements}
\addcontentsline{toc}{chapter}{Remerciements}
\markboth{REMERCIEMENTS}{REMERCIEMENTS}

Je tiens à exprimer ma reconnaissance à \textbf{M. Abdelmouttalib MAQIL}, mon
encadrant au sein d'EURAFRIC Information, pour la qualité de son suivi tout au
long de ces six semaines. Ses exigences — rattacher chaque composant au problème
qu'il traite, livrer une chaîne complète avant de la raffiner, et savoir défendre
chaque décision de conception sans notes — ont façonné ce travail bien au-delà de
sa dimension technique. Sa manière d'accueillir un résultat négatif comme un
résultat à part entière est, rétrospectivement, ce qui a rendu ce projet
intéressant.

Je remercie \textbf{EURAFRIC Information} de m'avoir accueilli et d'avoir proposé
un sujet à la fois ouvert et exigeant, ainsi que l'ensemble des collaborateurs
qui ont pris le temps de répondre à nos questions.

Mes remerciements vont également au corps professoral et à l'administration de
l'\textbf{Institut National des Postes et Télécommunications}, et en particulier
à la filière \textbf{Smart~ICT}, pour la formation reçue et pour le cadre donné à
ce projet de fin d'année.

Je remercie enfin mes camarades \textbf{Zakarya EL~WALI} et \textbf{Yasmine
BOUAJINE}, avec qui ce projet a été mené, pour la qualité des échanges et pour la
rigueur qu'ils ont apportée aux relectures — plusieurs des vérifications décrites
dans ce rapport sont nées d'une objection formulée par l'un ou l'autre.

Que ma famille et mes proches trouvent ici l'expression de ma gratitude pour leur
soutien constant.
